% Abstract (O-RAN / BMPP-DQN track).
%
% Unnumbered, per standard thesis convention; \addcontentsline gives it a
% Table of Contents entry despite \chapter* not numbering or listing it
% automatically. Placed in main_oran.tex between \maketitle and
% \tableofcontents, also the standard position.

\chapter*{Abstract}
\addcontentsline{toc}{chapter}{Abstract}

Open Radio Access Network (O-RAN) disaggregation exposes energy-relevant
control decisions -- Radio Unit (RU) activation, functional split
selection, transmit power, and physical resource block (PRB) allocation --
that are jointly discrete and continuous and that naturally span two
control timescales (Non-Real-Time and Near-Real-Time RIC). Existing deep
reinforcement learning approaches to this problem handle either the
discrete or the continuous side cleanly, or resolve the hybrid action
space by thresholding a continuous output, which destroys the discrete
decision's own gradient signal. This thesis proposes BMPP-DQN (Branching
Multi-Pass Parameterized DQN), a single coupled architecture that combines
branching discrete-action decomposition with multi-pass parameterized
continuous-action processing, to handle this hybrid, multi-timescale
action space directly.

BMPP-DQN is formulated as a parameterized Markov Decision Process over
four action branches -- RU activation, split selection, transmit power, and
PRB share -- and evaluated in a focused, single-gNB O-RAN simulation
environment against three baselines representing the discrete-only,
continuous-only, and existing-parameterized-action alternatives (DQN,
DDPG, and MP-DQN respectively), under a matched 3-seed, 500-episode
training and evaluation protocol. The comparison is cross-checked with
three independent multi-criteria decision-analysis schemes
(equal-weighted and entropy-weighted TOPSIS, and VIKOR), an
order-of-magnitude sensitivity sweep across the power model's unvalidated
constants, and a supplementary 10-seed statistical-power validation.

The results do not support this thesis's original target of a
$\geq15\%$ energy saving relative to the strongest baseline: BMPP-DQN's
mean power consumption is statistically indistinguishable from the
best-performing baseline's. What the evidence does support is a
multi-criteria balance across power, QoS satisfaction, and switching
stability that BMPP-DQN achieves more consistently than any single
baseline, a finding that holds across every power-model perturbation
tested and, substantially though not completely, across the larger
10-seed sample. A direct ablation further shows that this thesis's
two-timescale design -- holding discrete RU/split decisions fixed across
multiple environment steps rather than re-deciding them every step --
reduces switching frequency more than sixfold without a measurable cost to
reward, power, or QoS, isolating decision stability, not raw energy
efficiency, as the two-timescale design's actual contribution. These
findings are reported alongside the original energy-saving target rather
than in place of it, and the thesis's remaining limitations -- an
unvalidated power model and reward weighting, a single-gNB scope, and a
three-baseline comparison -- are stated explicitly as directions for
future work.
